% \section{Limitations}

% The study covers two domains, drug--disease facts with MONDO and species
% interactions with a taxonomy; recovery is evaluated on the first only. Questions are
% multiple choice with a fixed set of options, so abstention is measured as
% choosing ``Not enough information'' rather than as free-text refusal. The models are open and at most 14B, and the
% counterfactual edit covers only named questions whose answer has a MONDO parent.
% Recovery is evaluated with one model, and
% the page-built list uses exact category-name matching, which a learned family
% classifier could improve. The PrimeKG family-only group is small for
% contraindications (106 questions).

% \section{Conclusion}

% Models meet resolvable insufficiency with \clade{}: given a page that names only
% the answer's family, they stop declining and answer far less reliably, in
% medicine and in ecology. An ontology lists the answer once the family is known;
% finding the family and choosing within it remain open.


\section{Limitations}

Our study has several limitations. First, we evaluate two domains—biomedical drug--disease relations using MONDO and ecological species interactions using a biological taxonomy—and we run our ontology recovery pipeline only on the biomedical benchmark. Second, our abstention experiments use a multiple-choice format with an explicit \emph{Not enough information to tell} option rather than open-ended free-text generation, which allows exact control over distractor granularity and probability calibration but may differ from free-form refusal behavior. Third, we test open-weight models up to 14B parameters, and our counterfactual edit experiment is limited to named facts whose target entity has a parent category in the hierarchy. Fourth, our automated page-built recovery baseline relies on exact string matching to find category names in text; a trained relation-extraction model should substantially reduce the false-positive category matches that bottleneck Step~2. Finally, the number of available family-only contraindication facts in PrimeKG is modest (106 questions).

\section{Conclusion}

Retrieved text often describes facts at a coarser category level than a user's question demands. We show that this \emph{resolvable insufficiency} reliably triggers \clade{}: across both medicine and ecology, language models treat a coarse family or genus mention as enough reason to stop abstaining and guess a specific answer, cutting precision roughly in half. While a domain ontology can narrow the search to a short candidate list that holds the true answer 98.5\% of the time once the target family is isolated, extracting the right family from noisy pages and improving upon a simple popularity prior when picking the final member remain key open challenges for neuro-symbolic RAG.